\begin{frame}{확인 문제: 사전확률만으로 뒤집히는 메일}
  \begin{itemize}
    \item \textbf{(a)} 같은 학습 데이터에서, 사전확률
      \textbf{50/50}과 \textbf{90/10}(스팸/정상)일 때의
      $P(\text{spam}\mid\{$회의$\})$ 를 라플라스 스무딩으로 계산 ---
      어느 쪽에서 `회의'만 있는 메일을 \textbf{스팸}으로 분류하게 되는가?
      그 차이가 사전확률 $\log\frac{0.9}{0.1}\approx2.2$ nat에서
      오는 것인지, 가능도에서 오는 것인지 \textbf{구별}.
      {\footnotesize 답: 50/50 $\to$ 0.40 (정상), 90/10 $\to$
      $\frac{0.9\times0.5}{0.9\times0.5+0.1\times0.75}\approx0.857$
      (스팸). 가능도비 2/3은 동일, 바뀐 건 \textbf{사전확률 우위
      9:1}.}
    \item \textbf{(b)} 가능도비 $r$ 로 쓸 때: 50/50에서 정상 예측
      $\Leftrightarrow r<1$, 90/10에서 스팸 예측 $\Leftrightarrow
      9r>1 \Leftrightarrow r>1/9$.
      \vskip0.2em
      \textbf{$1/9 < r < 1$}인 메일만 50/50 $\to$ 90/10 변화로
      뒤집힘. 예: $r{=}0.5$ 면 1/3(정상) $\to$ 0.818(스팸).
      반대로, $r{=}2$(스팸 쪽)인 메일은 50/50 $\to$ 10/90 방향에서만
      뒤집힘. 한 문장: \textbf{가능도비가 1 근처($\sim$1/9 $\sim$ 9)인
      ``동점이 가까운'' 메일만, 사전확률의 9:1 변화에 의해 뒤집힘}.
  \end{itemize}
\end{frame}
